% \begin{frame}{Актуальність тематики досліджень}
%     На сьогоднішній день шкідливе програмне забезпечення (ШПЗ)
%     є однією з найважливіших загроз безпеці інформації.
%     Аналіз і розроблення прикладних математичних моделей для імітації поширення ШПЗ
%     є важливим завданням.
% \end{frame}
% \begin{frame}{Актуальність тематики досліджень}
%     На сьогоднішній день колоризація зображень
%     є важливою задачею комп'ютерного зору,
%     що застосовується для відновлення історичних матеріалів,
%     творчих індустрій та допоміжних технологій.

%     Сучасні методи колоризації на основі згорткових нейронних мереж (CNN),
%     генеративно-змагальних мереж (GAN), трансформерів (Transformer)
%     та дифузійних моделей (Diffusion) досягли значних результатів,
%     проте забезпечення інтерактивної взаємодії
%     та багатоваріантної генерації залишається складною задачею.

%     Новітні селективні моделі простору станів (SSM) є перспективними,
%     однак практично не застосовувалися для колоризації,
%     а їх ефективність порівняно з іншими генеративними підходами
%     систематично не вивчалась.

%     Це зумовлює потребу в розробленні інтерактивного програмного забезпечення,
%     яке через уніфікований пакет моделей використовує SSM, CNN, GAN,
%     Transformer та Diffusion для колоризації
%     і надає інструменти для об'єктивного та суб'єктивного порівняння.
% \end{frame}

% \begin{frame}{Актуальність тематики досліджень}
%     Колоризація зображень широко застосовується для відновлення історії,
%     творчих індустрій та допоміжних технологій, що робить її важливою задачею.

%     \vspace{1ex}

%     Наявні рішення, такі як CNN, GAN, Transformer та Diffusion, досягли значних результатів,
%     проте забезпечення інтерактивної взаємодії та багатоваріантної генерації
%     залишається складною задачею, в той час як новітні SSM
%     для колоризації не досліджені й не порівнювались із ними.

%     \vspace{1ex}

%     Враховуючи це, було створено програмне забезпечення для звичайних користувачів,
%     яке дозволяє користувачам обрати модель, кількість варіантів та інтенсивність підказок,
%     для проведення суб'єктивного порівняння,
%     а також окремо створено власну модель на базі SSM і
%     проведено об'єктивне порівняння цієї та наявних моделей за різними метриками.
% \end{frame}

% \begin{frame}{Актуальність тематики досліджень}
%     \begin{itemize}
%         \item Колоризація зображень широко застосовується для відновлення історії та 
%         творчих індустрій, що робить її важливою задачею.
    
%         % \vspace{1ex}
    
%         \item Наявні рішення, зокрема згорткові нейронні мережі (CNN),
%         генеративно-змагальні мережі (GAN), трансформери (Transformer)
%         та дифузійні моделі (Diffusion), досягли значних результатів,
%         проте забезпечення інтерактивної взаємодії та багатоваріантності
%         залишається складною задачею. Водночас новітні селективні моделі
%         простору станів (SSM) для колоризації не досліджені
%         й систематично не порівнювались із зазначеними підходами.
    
%         % \vspace{1ex}
    
%         \item Усунення цих прогалин потребує створення програмного забезпечення,
%         яке надасть користувачам можливість обирати модель, кількість варіантів
%         і інтенсивність підказок, а також проведення об'єктивного порівняння
%         моделей за метриками й суб'єктивного -- шляхом опитування користувачів.
%     \end{itemize}
% \end{frame}

% \begin{frame}{Актуальність тематики досліджень}
%     \begin{itemize}
%         \item Колоризація зображень широко застосовується для відновлення історії та 
%             творчих індустрій, що робить її важливою задачею.    

%         % \vspace{1ex}
    
%         \item Наявні рішення, зокрема згорткові нейронні мережі (CNN),
%             генеративно-змагальні мережі (GAN), трансформери (Transformer)
%             та дифузійні моделі (Diffusion), досягли значних результатів,
%             проте забезпечення інтерактивної взаємодії та багатоваріантності
%             залишається складною задачею. Водночас новітні селективні моделі
%             простору станів (SSM) для колоризації не досліджені
%             й систематично не порівнювались із зазначеними підходами.
    
%         % \vspace{1ex}
    
%         \item Усунення цих прогалин потребує створення програмного забезпечення,
%             яке надасть користувачам можливість обирати модель, кількість варіантів
%             і інтенсивність підказок, а також проведення об'єктивного порівняння
%             моделей за метриками й суб'єктивного -- шляхом опитування користувачів.
%     \end{itemize}
% \end{frame}
            % \item Сучасні архітектури (CNN, GAN, Transformer, Diffusion) зазвичай не забезпечують зручної \textbf{інтерактивності} та особливо \textbf{багатоваріантності}.

    
    % \begin{exampleblock}{Необхідність дослідження}
    %     Створення програмного забезпечення для інтерактивної колоризації з багатоваріантністю та проведення всебічного порівняльного аналізу архітектур.
    % \end{exampleblock}

\begin{frame}{Актуальність тематики досліджень}
    \begin{block}{Сучасний стан та проблематика}
        \begin{itemize}
            \item Існує високий попит на колоризацію у творчих індустріях та реставрації.
            \item Традиційні архітектури (CNN, трансформери) є \textbf{детермінованими},
                а генеративні архітектури (GAN, дифузійні моделі, авторегресивні трансформери),
                хоча й забезпечують \textbf{багатоваріантність},
                залишаються вкрай \textbf{ресурсоємними},
                що робить їх важкозастосовними для миттєвої \textbf{інтерактивної взаємодії у реальному часі}.
            \item Можливості селективних моделей простору станів (\textbf{SSM}) у цій сфері майже не вивчались.
        \end{itemize}
    \end{block}
\end{frame}